% !TEX root = ../notes.tex

\documentclass[../notes.tex]{subfiles}

\begin{document}

\section{March 18}
Here we go.

\subsection{The Localizations \texorpdfstring{$L_n$}{ Ln}}
Let's explain why our localizations $L_n^f$ may be difficult to understand.
\begin{example}
	Take $p>3$ so that $\mathbb S/(p,v_1)$ exists and is a ring, so we see that
	\[L_2^f\left(\frac{\mathbb S}{p,v_1}\right)=\frac{\mathbb S}{p,v_1}[1/v_2],\]
	Because $\mathbb S/(p,v_1)$ is bounded below (in fact, connective), we can try to compute its Adams--Novikov spectral sequence. One could hope to just take the localization of $v_2$ afterward to compute the above homotopy groups, but this program does not really work. Indeed, even if we have a class $x_0$ which has $v_2^3x=0$ (so that it vanishes after inverting $v_2$), we only know that $x_0$ to a class $\widetilde x_0$ which has $v_2^3\widetilde x_0=\widetilde x_1$ for some other class $\widetilde x_1$ in a higher filtration. In fact, it turns out that $\widetilde x_1$ looks like a $p$th power, so we will just write it as $\widetilde x_1^p$. In fact, this sort of thing can go on infinitely, finding more and more $p$th roots in these filtration jumps.
\end{example}
However, there are variant functors to the $L_n^f$s which are easier to compute.
\begin{definition}
	We define $L_n$ as the localization $L_{K(0)\oplus\cdots\oplus K(n)}$.
\end{definition}
\begin{remark}
	By construction, a spectrum $M$ is $L_n$-acyclic if and only if $M\otimes(K(0)\oplus\cdots\oplus K(n))=0$. Thus, $M$ is $L_n^f$-acyclic if and only if $M$ is a colimit of finite $L_n$-acyclic spectra; it turns out that $L_n$ is smashing, but this is far from obvious. It will turn out that every nonzero idempotent algebra $R$ sits in the sequence
	\[L_n^f\mathbb S\to R\to L_n\mathbb S\]
	for some $n\ge0$. (The map $L_n^f\mathbb S\to L_n\mathbb S$ exists formally by the functoriality of the localization.)
\end{remark}
\begin{remark}
	It turns out that $L_0=L_0^f$ and $L_1=L_1^f$, but $L_n\ne L_n^f$ for all $n\ge2$.
\end{remark}
These localizations $L_n$ are still interesting to us because they are able to recover $X$.
\begin{theorem}[Hopkins--Ravenel] \label{thm:recover-x-from-localizations}
	Fix a perfect $p$-local spectrum $X$. Then
	\[X=\lim L_\bullet X.\]
\end{theorem}
\begin{remark}
	It is currently open if $X$ is the limit of the $L_n^fX$s.
\end{remark}
Now, if $X$ is a $p$-local spectrum, then $L_{K(n)}X$ can be understood with Galois theory; we will talk more about this after spring break. For example, we will be able to prove the following.
\begin{theorem}[Hopkins--Ravenel] \label{thm:galois-theory}
	If $X$ is a finite spectrum of type $n$, then $L_{K(n)}X$ is in the thick subcategory of $p$-local spectra generated by $K(n)$.
\end{theorem}
\begin{remark}
	For example, it follows that $X$ has finite-dimensional homotopy groups in each degree because this is true for $K(n)$.
\end{remark}
\begin{example}
	Fix an odd prime $p$. We will be able to show that $L_{K(1)}\mathbb S/p$ is the fiber of some map $K(1)\to K(1)$; this implies the result for $n=1$ (in this case) by \Cref{thm:thick-finite-subcat}.
\end{example}

\subsection{The Smashing Theorem}
Now that we are motivated to understand these localizations $L_{K(n)}$, let's show that $L_n$ is smashing.
\begin{theorem} \label{thm:ln-smashing}
	Fix a nonnegative integer $n$. The localization $L_n$ is smashing.
\end{theorem}
\begin{proof}
	One checks by hand that $L_0=L_\QQ$ is smashing. Let's next prove that $L_1$ is smashing. To this end, we need to show that $L_1\mathbb S\otimes X$ is $L_1$-local for any $X$. Given this, one knows that the canonical map $\mathbb S\otimes X\to L_1\mathbb S\otimes X$ is an equivalence after tensor product with $K(0)\otimes K(1)$, so we are able to recognize our localization as $L_1\mathbb S\otimes-$, which is then manifestly smashing.

	To this end, note that we have a pullback square
	% https://q.uiver.app/#q=WzAsNCxbMCwwLCJMXzFcXG1hdGhiYiBTXFxvdGltZXMgWCJdLFsxLDAsIihMXzFcXG1hdGhiYiBTXFxvdGltZXMgWCleXFxsYW5kX3AiXSxbMCwxLCJcXFFRXFxvdGltZXMgTF8xXFxtYXRoYmIgU1xcb3RpbWVzIFgiXSxbMSwxLCJcXFFRXFxvdGltZXMoTF8xXFxtYXRoYmIgU1xcb3RpbWVzIFgpXlxcbGFuZF9wIl0sWzAsMV0sWzEsM10sWzIsM10sWzAsMl1d&macro_url=https%3A%2F%2Fraw.githubusercontent.com%2FdFoiler%2Fnotes%2Fmaster%2Fnir.tex
	\[\begin{tikzcd}[cramped]
		{L_1\mathbb S\otimes X} & {(L_1\mathbb S\otimes X)^\land_p} \\
		{\QQ\otimes L_1\mathbb S\otimes X} & {\QQ\otimes(L_1\mathbb S\otimes X)^\land_p}
		\arrow[from=1-1, to=1-2]
		\arrow[from=1-1, to=2-1]
		\arrow[from=1-2, to=2-2]
		\arrow[from=2-1, to=2-2]
	\end{tikzcd}\]
	which is induced from some general nonsense. Indeed, for any $R$-modules $A$ and $B$ for which $L_BL_A=0$, there is a pullback square
	% https://q.uiver.app/#q=WzAsNCxbMCwwLCJMX3tBXFxvcGx1cyBCfU0iXSxbMSwwLCJMX0JNIl0sWzEsMSwiTF9BTF9CTSJdLFswLDEsIkxfQU0iXSxbMCwzXSxbMywyXSxbMCwxXSxbMSwyXV0=&macro_url=https%3A%2F%2Fraw.githubusercontent.com%2FdFoiler%2Fnotes%2Fmaster%2Fnir.tex
	\[\begin{tikzcd}[cramped]
		{L_{A\oplus B}M} & {L_BM} \\
		{L_AM} & {L_AL_BM}
		\arrow[from=1-1, to=1-2]
		\arrow[from=1-1, to=2-1]
		\arrow[from=1-2, to=2-2]
		\arrow[from=2-1, to=2-2]
	\end{tikzcd}\]
	from which we recover the earlier pullback by taking $R=\mathbb S_{(p)}$ and $A=\QQ$ (so that $L_AM=M\otimes\QQ$) and $B=\mathbb S/p$ (so that $L_BM=\lim M/p^\bullet=M^\land_p$). For example, one has that $L_{\QQ\oplus\mathbb S/p}M=M$ because $(\QQ\oplus\mathbb S/p)\otimes M=0$ implies that $M=0$.
	
	Now, pullbacks of local objects are local, so it is enough to show that the top-right and bottom-left corners of the square are local. For example, any $\QQ$-module is $L_0$-local because $L_0=\QQ\otimes-$ by the induction. Because $L_1=L_{K(0)}\oplus L_{K(1)}$, we thus see that the bottom-left corner is $L_1$-local as well. As for the top-right corner, we see that
	\[(L_1\mathbb S\otimes X)^\land_p=\lim (L_1\mathbb S\otimes X)/p^\bullet.\]
	Because being local is preserved under limits, we may merely show that each individual module in the limit is local. By taking extensions, it is enough to just show that $(L_1\mathbb S\otimes X)/p$ is local.
	
	To this end, we claim that
	\[L_1\mathbb S/p\stackrel?=L_{K(1)}\mathbb S/p.\]
	Indeed, there is certainly a map $L_1\mathbb S\to L_{K(1)}\mathbb S$ by taking the localization. It is an equivalence after $K(1)_*$ due to the localization, and it is also an equivalence after taking $K(0)_*$ modulo $p$ because $K(0)$ is just $\QQ$.

	Now, $(L_1\mathbb S\otimes X)/p$ is
	\[L_1\mathbb S/p\otimes X=L_{K(1)}\mathbb S/p\otimes X\]
	by the claim. But Galois theory in the form of \Cref{thm:galois-theory} implies that $L_{K(1)}\mathbb S/p\otimes X$ is in the thick subcategory generated by $K(1)\otimes X$, which is local!

	Let's briefly explain how we may continue the induction. We would like to show that $L_2$ is smashing, so we need to show that $L_2\mathbb S\otimes X$ is $L_2$-local. One starts by showing that $(L_2\mathbb S\otimes X)/(p^{a_0},v_1^{a_1})$ is $L_2$-local for some large $a_0$ and $a_1$, and then we show that $(L_2\mathbb S\otimes X)/p^{a_0}$ is $L_2$-local, and then we get $L_2\mathbb S\otimes X$ being $L_2$-local. To even get started at the first step, one shows that $L_2\mathbb S/(p^{a_0},v_1^{a_1})$ is $L_{K(2)}\mathbb S/(p^{a_0},v_1^{a_1})$ for some $a_0$ and $a_1$ as above. One then uses pullback squares and limiting arguments to make the reductions.
\end{proof}
The localizations also do not communicate with each other too much.
\begin{corollary}
	Fix a $p$-local spectrum $X$. Then $L_{K(n)}L_{n-1}X=0$.
\end{corollary}
\begin{proof}
	We would like to show that the map $L_{n-1}X\to0$ is the $K(n)$-localization map. Well, $0$ is certainly local, so it remains to show that we have an equivalence after $K(n)\otimes-$, so we need to show that $K(n)\otimes L_{n-1}X=0$. By \Cref{thm:ln-smashing}, this is $K(n)\otimes L_{n-1}\mathbb S\otimes X=0$, so it is enough to show that $L_{n-1}K(n)=0$. But this amounts to checking that $(K(0)\oplus\cdots\oplus K(n-1))\otimes K(n)=0$, which follows by the distinctness of our characteristics.
\end{proof}
\begin{remark}
	The formal nonsense from earlier produces the following pullback square.
	% https://q.uiver.app/#q=WzAsNCxbMCwwLCJMX25YIl0sWzEsMCwiTF97SyhuKX1YIl0sWzAsMSwiTF97bi0xfVgiXSxbMSwxLCJMX3tuLTF9TF97SyhuKX1YIl0sWzAsMV0sWzEsM10sWzIsM10sWzAsMl1d&macro_url=https%3A%2F%2Fraw.githubusercontent.com%2FdFoiler%2Fnotes%2Fmaster%2Fnir.tex
	\[\begin{tikzcd}[cramped]
		{L_nX} & {L_{K(n)}X} \\
		{L_{n-1}X} & {L_{n-1}L_{K(n)}X}
		\arrow[from=1-1, to=1-2]
		\arrow[from=1-1, to=2-1]
		\arrow[from=1-2, to=2-2]
		\arrow[from=2-1, to=2-2]
	\end{tikzcd}\]
	Thus, the difference between $L_nX$ and $L_{n-1}X$ is more or less given by $L_{K(n)}X$. For example, $\lim L_\bullet X$ (which agrees with $X$ if $X$ is a finite spectrum) is a limit of $L_{K(a_1)}\cdots L_{K(a_m)}X$s, where $a_1<\cdots<a_m$ is an increasing sequence. Such iterated localizations is understood using Galois theory.
\end{remark}
We now have the tools to state the Chromatic splitting conjecture.
\begin{conj}[Chromatic spliting] \label{conj:chromatic-splitting-again}
	Fix a positive integer $n\ge1$. If $X$ is a finite $p$-local spectrum, then the map
	\[L_{n-1}X\to L_{n-1}L_{K(n)}X\]
	is the inclusion of a direct summand.
\end{conj}
\begin{remark}
	This is known at all primes for $n\in\{1,2\}$.
\end{remark}
In the ``large $p$'' regime, one can reduce the conjecture to a question in pure algebra about the cohomology of some perfectoid spaces.
\begin{remark}
	Assuming \Cref{conj:chromatic-splitting-again}, one finds that $\mathbb S_p^\land$ is a direct summand of $\prod_{n\ge1}L_{K(n)}\mathbb S$. Thus, one could understand the homotopy groups of $\mathbb S_p^\land$ by only looking at the completions and trying to glue them together.
\end{remark}

\end{document}